\documentclass[12pt]{article}
\usepackage[T1]{fontenc}
\usepackage[utf8]{inputenc}
\usepackage{lmodern}
\usepackage{textcomp}
\usepackage{enumitem}
\usepackage{geometry}
\geometry{margin=1in}
\begin{document}
\begin{center}
Answer Key - Religious Studies - K-12 Level Assessment 18 \
\small (Answers are ordered exactly as in the question paper.)
\end{center}

\section*{Multiple Choice}
\begin{enumerate}[label=\arabic*.]
\item \textbf{Answer:} B
\\ \textit{Options:} A) 1947; B) 1948; C) 1945; D) 1949
\\ \textit{Note:} The State of Israel was established on May 14, 1948, following the end of the British Mandate for Palestine and the approval of the United Nations Partition Plan for Palestine in 1947, which aimed to create separate Jewish and Arab states.
\item \textbf{Answer:} D
\\ \textit{Options:} A) Zazen; B) Tariki; C) Kami; D) Ikebana
\\ \textit{Note:} Ikebana is the traditional Japanese art of flower arranging, characterized by its emphasis on harmony, balance, and the aesthetic use of space, making it the correct answer to the question.
\item \textbf{Answer:} B
\\ \textit{Options:} A) Ritual texts; B) Philosophical texts; C) Hymns; D) Origin stories
\\ \textit{Note:} The Upanishads are characterized as philosophical texts because they explore fundamental questions about existence, the nature of reality, and the self, which are central themes in Indian philosophy and spirituality.
\item \textbf{Answer:} C
\\ \textit{Options:} A) Paul of Tarsus; B) Clare of Assisi; C) Francis of Assisi; D) John of the Cross
\\ \textit{Note:} Francis of Assisi is known for his deep connection with nature, which is exemplified by his preaching to birds, and he is also famously associated with stigmata, marking him as a significant figure in Christian mysticism.
\item \textbf{Answer:} C
\\ \textit{Options:} A) 1749-1945; B) 1052-1616; C) 1641-1853; D) 1517-1870
\\ \textit{Note:} The timespan of 1641-1853 marks the period during which Japan implemented a policy of sakoku, or national isolation, which included the expulsion of Christians and the restriction of foreign influence, particularly from Europe.
\end{enumerate}

\section*{True / False}
\begin{enumerate}[label=\arabic*.]
\setcounter{enumi}{5}
\item \textbf{Answer:} False
\\ \textit{Options:} A) True; B) False
\\ \textit{Note:} The statement is false because Confucian 'sage kings' are idealized figures who embody moral virtue and wisdom but are characterized by their passive governance, emphasizing the importance of moral example over direct political engagement.
\item \textbf{Answer:} False
\\ \textit{Options:} A) True; B) False
\\ \textit{Note:} Guru Ram Das was not martyred for rejecting Islam; instead, he was the fourth Sikh Guru known for his contributions to Sikhism and promoting love, equality, and community service.
\item \textbf{Answer:} False
\\ \textit{Options:} A) True; B) False
\\ \textit{Note:} The statement is false because while the Qur'an is a central religious text in Islam, it is not universally regarded as the sole symbol of the religion, as many Muslims may consider other elements, such as the Prophet Muhammad or the mosque, to also represent their faith.
\item \textbf{Answer:} False
\\ \textit{Options:} A) True; B) False
\\ \textit{Note:} In Jainism, liberation is sought from both Punya (good karma) and Papa (bad karma), as both types of karma bind the soul to the cycle of rebirth, indicating that liberation involves transcending all forms of karma rather than solely seeking freedom from good karma.
\item \textbf{Answer:} False
\\ \textit{Options:} A) True; B) False
\\ \textit{Note:} The statement is false because the Ka'ba, located in Mecca, is primarily associated with the Islamic tradition and is not historically linked to prophets like Moses, who lived in a different geographic and cultural context.
\end{enumerate}

\section*{Long Form Response}
\begin{enumerate}[label=\arabic*.]
\setcounter{enumi}{10}
\item \textbf{Answer:} The Hizmet movement, led by Fetullah Gülen, originated in Turkey in the 1970 s and is rooted in Islamic teachings that emphasize education, dialogue, and social service. Its core beliefs focus on the importance of interfaith dialogue, modern education, and community service, aiming to promote a peaceful, tolerant society. The movement has significantly impacted Turkish society by establishing schools, media outlets, and charitable organizations, thereby influencing education and public discourse. However, it has also been involved in political controversies, particularly after a fallout with the Turkish government, which accused it of attempting to undermine state authority. Overall, the Hizmet movement highlights the intersection of religion and social activism in contemporary Turkey.
\\ \textit{Note:} correct because it accurately outlines the origins of the Hizmet movement in the 1970 s, its foundational beliefs centered on education and interfaith dialogue, and its substantial societal impact through the establishment of educational institutions and media, reflecting its goals of fostering tolerance and community service.
\item \textbf{Answer:} The Caliphs played a crucial role as both religious and administrative leaders in the early Islamic community after the death of the Prophet Muhammad. They were considered the successors to the Prophet and were responsible for guiding the Muslim community in spiritual matters and upholding Islamic law. The first four Caliphs, known as the "Rightly Guided Caliphs," helped to expand the Islamic state, maintain unity among Muslims, and establish essential practices and governance. Their leadership was significant in shaping the early Islamic faith and community, as they set precedents for future governance and interpretation of Islamic teachings. Thus, the Caliphs were vital in fostering a sense of identity and cohesion within the rapidly growing Muslim society.
\\ \textit{Note:} correct because it accurately identifies the Caliphs as both spiritual successors to the Prophet Muhammad and pivotal administrators who fostered community unity, expanded Islamic influence, and established foundational practices in governance and law.
\end{enumerate}

\vfill
\noindent\textit{End of Answer Key}
\end{document}